\documentclass[11pt]{article}
\usepackage[T1]{fontenc}
\usepackage{lmodern}
\usepackage[margin=1in]{geometry}
\usepackage{amsmath,amssymb,amsthm}
\usepackage[hidelinks]{hyperref}
\hypersetup{pdftitle={Xu--Tang Proof Audit},pdfauthor={Christopher D. Long}}
\newtheorem{lemma}{Lemma}
\newcommand{\Pbb}{\mathbb P}
\newcommand{\Ebb}{\mathbb E}
\title{Xu--Tang Proof Audit}
\author{Christopher D. Long}
\date{}
\begin{document}
\maketitle
\begin{abstract}
This note explains why the earlier moving-tail and terminal-cloud proof
assemblies were replaced by a fixed-time fluid limit and a simultaneous
gamma-clock comparison. It distinguishes proof gaps from false conclusions
and repairs a reusable martingale inequality. The companion paper,
\emph{Extremes and Moments in Greedy Coupon Collection}, contains the complete
replacement proof for the fixed-parameter completion results. The old
moving-tail fluctuation program is not an input to that proof.
\end{abstract}

\section{Scope and lineage}
The local moving-tail synthesis explicitly retained a common-interval
martingale hypothesis. The local terminal-cloud synthesis claimed an
unconditional completion law but retained a uniform-integrability hypothesis
for the sharp mean. Their four component proofs were substantially duplicated.
Earlier research discussions described later revisions: an integrating-factor
envelope proof, localized martingale control, and a hybrid lower-layer
remainder. Those discussions establish development history, not the validity
of unavailable source proofs. The findings below concern the inspected local
files; they do not assert that every later revision retained the same errors.

The new paper supersedes both proof assemblies. It retains their exact
histogram drift, deterministic amplitude, and target completion law, and
supplies a new proof with stronger conclusions. The original local files are
preserved outside this repository, with fingerprints in the audit records.
The public sources use descriptive names, author attribution, no dynamic date
or submission contact block, and one canonical proof chain. The broader
terminal-defect framework remains a separate manuscript.

\section{Intervals, stopping rules, and quantifiers}
Write $c=m-1$, $H_t=gt/n$, and
\[
 A_j(t)=ne^{-H_t}(1+H_t)^j,\qquad W_j(t)=\sqrt{A_j(t)\log n}+1.
\]
The fixed-output martingale component controls row $j$ only to
\[
 T^-_{j,n}=\frac ng\{\log n+(j-1)\log\log n+B-\omega_n\}+O(1).
\]
The all-row bootstrap needs every row $j\le c$ controlled to $T^-_{c,n}$.
For $j<c$ these intervals differ. An estimate on the shorter interval is
not a proof on the longer one. Likewise, stopping when any row exits an
envelope differs from stopping only at the top-row exit. Variance estimates
must be justified for the stopping rule actually used. Later discussions
identified this and proposed a localized replacement.

Also $X_n=O_{\Pbb}(1)$ means that for every $\varepsilon>0$ there is a
finite $C_\varepsilon$ with
$\limsup_n\Pbb(|X_n|>C_\varepsilon)\le\varepsilon$.
It does not imply one finite $C$ with $\Pbb(|X_n|\le C)\to1$.
A bootstrap can retain the former quantifiers, choosing its height after
$\varepsilon$, or establish a stronger launch estimate. The replacement
paper uses an explicit exponentially small failure probability at fixed time.

A sign convention also needs repair. If
$D_j(t+1)=\mathbf1\{L_j(t+1)=L_j(t)-1\}$ and $\xi_j=D_j-p_j$,
the noise in the recursion for $L_j$ is $-\xi_j$. Alternatively define
$\zeta_j=p_j-D_j$ and use a plus sign. Absolute-value estimates survive,
but representation identities must use the same convention throughout.

\section{What the envelope does not prove}
The component envelope has a cushion:
\[
 L_j(t)\le C\{ne^{-H_t}(1+H_t)^{j+\delta}+\log n\},\qquad\delta>0.
\]
Omitting $(1+H_t)^\delta$ strengthens it without proof. Earlier revisions
explicitly repaired this mismatch.

Even an uncushioned density-scale envelope is too coarse to directly control
the centered quadratic remainder at the square-root scale. For $c=1$,
$g\ge2$, and $H_{\tau_n}=\eta\log n+O(1/n)$, $0<\eta<1/9$,
it permits $|Z_0(s)|$ of order $n^{1-\eta}$ on
$s=\tau_n+O(n)$. A diagonal propagator has order one there. The resulting
available quadratic bound in the top equation can be of size
\[
 \frac1{n^2}\sum_{s=\tau_n}^{\tau_n+\Theta(n)}
 \Gamma_{11}(s+1,t)Z_0(s)^2\asymp n^{1-2\eta}.
\]
Since $W_1(t)\asymp n^{(1-\eta)/2}\log n$, its ratio to the target scale is
\[
 n^{(1-3\eta)/2}/\log n\longrightarrow\infty.
\]
This invalidates an inference from the envelope; it does not disprove the
actual fluctuation theorem. The later hybrid proposal was intended to add
the missing early control of centered lower layers. The replacement proof
needs no nonlinear convolution estimate.

The original upper-envelope crossing proof has a separate issue. Its exponent
has the form
\[
 \frac{(R/4+S)^2}{C(V+R+S)},\qquad R=A\log n,
\]
with drift surplus approximately $S\asymp V/H$ and $H\asymp\log n$.
The allowed scale $V\asymp RH$, $S\asymp R$ yields order $R/H=O(A)$,
not $A\log n$. Chopping into subexcursions does not automatically renew the
buffer. A valid integrating-factor or renewed-buffer argument is needed.
Later discussions report an integrating-factor repair; the new paper uses
neither version of this envelope lemma.

\section{Probability convergence does not control moments}
A lower-layer survivor count that is $o_{\Pbb}(1)$ need not have vanishing
factorial moments. For example,
\[
 X_n=n\text{ with probability }n^{-2},\qquad X_n=0\text{ otherwise}
\]
converges in probability to zero, while
$\Ebb[X_n(X_n-1)]\to1$. Likewise, conditional factorial-moment estimates
on events of probability $1-o(1)$ need not give unconditional moment
convergence: exceptional events may carry the moments.

Bounded conditional generating functions or Laplace tests are enough for
weak convergence. Moment conclusions require actual tail or $L^p$ control.
The new proof supplies both: independent-gamma Laplace functionals, explicit
extreme-value tail bounds, and an $L^p$ clock error. Exponentially unlikely
fixed-time failures are handled with coarse global moments.

A bounded next factorial moment can justify uniform integrability under
conditional laws, when that bound holds in the required conditional sense.
It does not by itself control the complement of the good event.

The base case $c=0$ also needs consistent entrance bookkeeping. An entrance
at $H_s=\log n+y-\omega_n$ leaves a window $\omega_n$ and mass of order
$e^{\omega_n}$. An entrance using the common $(c-1)\log\log n$ formula
leaves an extra $\log\log n$ window and a mass of order
$e^{\omega_n}\log n$. Either convention works if used consistently.
The gamma proof handles $m=1$ directly.

Finally, $\{0\}$ is not a continuity set for a Poisson law on $\mathbb R$.
For integer-valued variables use $(-1/2,1/2)$ to extract the zero mass from
weak convergence, or work explicitly in the discrete topology.

\section{Predictable-variance truncation repaired}
Stopping a discrete martingale at the first time $V_t>b$ includes a possible
variance overshoot. It does not ensure a stopped bracket at most $b$.
The following predictable truncation repairs the launch component's argument.

\begin{lemma}\label{lem:variance}
Let $(M_t,\mathcal F_t)_{t=0}^N$ be a square-integrable martingale, $M_0=0$,
and let
\[
 V_t=\sum_{s=1}^t\Ebb[(M_s-M_{s-1})^2\mid\mathcal F_{s-1}].
\]
For $a,b>0$, $\Pbb(|M_N|>a)\le\Pbb(V_N>b)+b/a^2$.
\end{lemma}
\begin{proof}
Since $V_t$ is $\mathcal F_{t-1}$-measurable,
\[
 \widetilde M_N=\sum_{t=1}^N(M_t-M_{t-1})\mathbf1_{\{V_t\le b\}}
\]
is a martingale transform. Its bracket is
\[
 \sum_{t=1}^N\mathbf1_{\{V_t\le b\}}(V_t-V_{t-1})\le b,
\]
because $V_t$ is nondecreasing. Hence $\Ebb\widetilde M_N^2\le b$.
On $\{V_N\le b\}$, $\widetilde M_N=M_N$. Chebyshev's inequality and a
split over this event give the result. The overshooting increment is excluded.
\end{proof}

\section{Replacement proof and status}
The replacement dependencies are
\[
\begin{gathered}
 \text{exact class rates and fixed-time concentration}\\
 \Downarrow\\
 \text{sparse population and base-layer amplitude}\\
 \Downarrow\\
 \text{independent gamma clocks with a simultaneous }L^p\text{ error}\\
 \Downarrow\\
 \text{Poisson completion points, extreme laws, all fixed moments}\\
 \Downarrow\\
 \text{discrete-time results by independent clock comparison.}
\end{gathered}
\]
This proves the fixed-$m$, fixed-batch completion law, its sharp mean and
variance, and joint late order statistics. It extends to bounded random
batch laws, replacement sampling, and independent holding. The deterministic
layer asymptotics remain, with a short triangular ODE proof. A finite
sparse-state mean bound also gives a certified approximation for computation.

The separate all-layer moving-tail fluctuation estimate is neither assumed
nor claimed. A logarithmically buffered square-root bound is not a normal
limit theorem. Growing parameters and nonuniform offers remain further
directions. This is an analytical and computational audit, not a
proof-assistant verification or an independent referee report.
\end{document}
